% Abhakseite „Das kann ich" – jede Hauptnummer genau einmal, nach Zweigen gruppiert, ohne Lösungen
\clearpage
\hypertarget{abhaken}{}\noindent{\large\bfseries Das kann ich}\par\smallskip
\noindent{\small Hake am Ende ab, was du kannst.}\par
\begingroup\setlength{\parskip}{0pt}
\abhakgruppe{Kennst du schon}
\abhak{1}{Ich kann negative Zahlen addieren und subtrahieren.}
\abhak{2}{Ich kann mit Vorzeichen multiplizieren und Punkt vor Strich beachten.}
\abhak{3}{Ich kann den Wert eines Terms berechnen.}
\abhak{4}{Ich kann gleichartige Glieder zusammenfassen.}
\abhak{5}{Ich finde den Fehler beim Zusammenfassen.}
\abhak{6}{Ich kann Terme malnehmen.}
\abhak{7}{Ich kann Quadratzahlen bilden und erkennen.}
\abhakgruppe{Einheit 1 · Klammern auflösen}
\abhak{8}{Ich erkenne die Glieder eines Terms mit ihrem Vorzeichen.}
\abhak{9}{Ich erkenne, was vor einer Klammer steht.}
\abhak{10}{Ich kann eine Klammer mit Plus davor weglassen.}
\abhak{11}{Ich kann eine Zahl mit einer Klammer malnehmen.}
\abhak{12}{Ich kann eine Klammer mit Minus davor auflösen.}
\abhak{13}{Ich finde den Fehler beim Auflösen von Klammern.}
\abhak{14}{Ich kann begründen, ob zwei Terme gleichwertig sind.}
\abhak{15}{Ich kann eine Sachaufgabe mit einer Klammer lösen.}
\abhakgruppe{Einheit 2 · Ausklammern}
\abhak{16}{Ich kann einen gemeinsamen Faktor ausklammern.}
\abhak{17}{Ich finde den Fehler beim Ausklammern.}
\abhak{18}{Ich kann entscheiden, ob man einen Faktor ausklammern kann.}
\abhak{19}{Ich kann mit Ausklammern eine Seitenlänge bestimmen.}
\abhakgruppe{Einheit 3 · Summe mal Summe}
\abhak{20}{Ich erkenne, welches Glied mit welchem malgenommen wird.}
\abhak{21}{Ich kann den Wert eines Terms mit zwei Variablen berechnen.}
\abhak{22}{Ich kann Terme mit zwei Variablen zusammenfassen.}
\abhak{23}{Ich kann zwei Klammern miteinander malnehmen.}
\abhak{24}{Ich kann zwei Klammern miteinander malnehmen – weiter.}
\abhak{25}{Ich finde den Fehler beim Malnehmen zweier Klammern.}
\abhak{26}{Ich kann mit einem Rechteck begründen, warum vier Produkte entstehen.}
\abhak{27}{Ich kann mit Klammer mal Klammer Flächen vergleichen.}
\abhakgruppe{Einheit 4 · Binomische Formeln}
\abhak{28}{Ich erkenne, ob eine Klammer mit sich selbst malgenommen wird.}
\abhak{29}{Ich erkenne, was $a$ und was $b$ ist.}
\abhak{30}{Ich erkenne, welche binomische Formel passt.}
\abhak{31}{Ich kann eine Klammer mit einer binomischen Formel ausmultiplizieren.}
\abhak{32}{Ich kann eine Klammer mit einer binomischen Formel ausmultiplizieren – weiter.}
\abhak{33}{Ich kann eine Klammer mit einer binomischen Formel ausmultiplizieren – weiter.}
\abhak{34}{Ich kann mit einer binomischen Formel im Kopf rechnen.}
\abhak{35}{Ich finde den Fehler bei einer binomischen Formel.}
\abhak{36}{Ich finde den Fehler bei einer binomischen Formel – weiter.}
\abhak{37}{Ich kann begründen, warum $(a + b)^2$ nicht $a^2 + b^2$ ist.}
\abhak{38}{Ich kann mit einer binomischen Formel die Fläche eines Rahmens berechnen.}
\abhakgruppe{Einheit 5 · Faktorisieren}
\abhak{39}{Ich erkenne, ob ein Term ein Quadrat ist.}
\abhak{40}{Ich erkenne, ob das Mittelglied passt.}
\abhak{41}{Ich kann eine Summe mit einer binomischen Formel in ein Produkt verwandeln.}
\abhak{42}{Ich kann eine Summe mit einer binomischen Formel in ein Produkt verwandeln – weiter.}
\abhak{43}{Ich kann eine Summe mit einer binomischen Formel in ein Produkt verwandeln – weiter.}
\abhak{44}{Ich finde den Fehler beim Faktorisieren.}
\abhak{45}{Ich kann begründen, wann man mit einer binomischen Formel faktorisieren kann.}
\abhak{46}{Ich kann die Seitenlänge eines Quadrats aus seiner Fläche als Term bestimmen.}
\endgroup
